\chapter{Системийн зохиомж ба хөгжүүлэлт}
\label{Chapter3}

\section{Ерөнхий архитектур}
Хөгжүүлсэн систем нь урд талын хэсэг, арын хэсэг, өгөгдлийн сан, хиймэл оюуны уялдуулалт, нэвтрүүлэлтийн орчин гэсэн таван үндсэн давхаргатай. Архитектурын зорилго нь хэрэглэгчийн интерфейс, бизнес логик, өгөгдөл хадгалалт, хиймэл оюуны үйлчилгээ, орчны тохиргоог хооронд нь хэт нягт хамааралгүйгээр уялдуулах явдал юм.
\insertimage{zurag-diagram/system-architecture.png}{Системийн ерөнхий архитектурын диаграм}

Системийн ерөнхий өгөгдлийн урсгалыг дараах байдлаар тайлбарлаж болно.
\begin{enumerate}
\item Хэрэглэгч React дээрх урд талын хэсгээр дамжин API хүсэлт илгээнэ.
\item Django REST дээрх арын хэсэг хүсэлтийг хүлээн авч баталгаажуулалт болон эрх шалгана.
\item Шаардлагатай тохиолдолд PostgreSQL эсвэл SQLite өгөгдлийн сангаас өгөгдөл уншиж, шинэчилнэ.
\item Хиймэл оюунтай холбоотой төгсгөлийн цэгүүд хичээл эсвэл курсийн агуулгаас эх текст бэлтгэн Gemini API дуудах эсвэл нөөц логик ашиглана.
\item Хариу нь JSON хэлбэрээр урд талын хэсэгт очиж, хуудасны бүрэлдэхүүнүүд дээр дүрслэгдэнэ.
\end{enumerate}

Энэхүү архитектурын давуу тал нь урд болон арын хэсгийг тус тусад нь хөгжүүлэх, турших, өргөтгөх боломжтойд оршино. Мөн хиймэл оюуны үйлчилгээ системийн гадна оршиж байгаа тул тухайн давхаргын доголдлыг бусад үндсэн модулиудаас тодорхой хэмжээнд тусгаарлаж чадна.

\section{Технологийн стек сонголт}
\subsection{Хэрэглэгчийн тал}
Урд талын хэсэгт React, TypeScript, Vite ашигласан \cite{react,typescript,vite}. React нь бүрэлдэхүүнд суурилсан бүтэцтэй тул хуудас, бүтэц, жижиг хэрэгсэл, карт, цонх зэрэг интерфейсийн элементүүдийг дахин ашиглахад хялбар болгодог. TypeScript ашигласнаар \texttt{Lesson}, \texttt{Course}, \texttt{LessonProgress}, \texttt{CommunityPost} зэрэг өгөгдлийн төрлүүдийг тодорхойлж, API хариу болон интерфейсийн төлөвийн нийцтэй байдлыг сайжруулсан. Vite нь хурдан хөгжүүлэлтийн сервер, богино угсралтын хугацаа, орчин үеийн модуль системээр урд талын хөгжүүлэлтийн бүтээмжийг нэмэгдүүлсэн.

\subsection{Серверийн тал}
Арын хэсэгт Django болон Django REST Framework ашигласан \cite{django,drf}. Django ORM нь хэрэглэгч, курс, хичээл, ахиц, хамтын орчны өгөгдлийн холбоотой бүтэц дээр ажиллахад тохиромжтой. DRF ашигласнаар жагсаах, авах, үүсгэх API-уудыг ерөнхий view хэлбэрээр богино бөгөөд өргөтгөх боломжтой байдлаар хэрэгжүүлсэн. Баталгаажуулалтад \texttt{djangorestframework-simplejwt} ашигласан нь нэг хуудсанд суурилсан урд талын бүтэцтэй нийцтэй шийдэл болсон.

\subsection{Өгөгдөл хадгалалт ба орчны тохиргоо}
Хөгжүүлэлтийн үеийн тохиргоонд SQLite ашиглах боломжтой, харин Docker орчинд PostgreSQL ашиглаж байгаа \cite{postgres}. Энэ нь хөгжүүлэгч дангаараа хурдан турших болон багийн орчинд тогтвортой ажиллуулах хоёр хэрэгцээг зэрэг хангаж байна. Арын хэсгийн тохиргоонд \texttt{DJANGO\_DEBUG}, \texttt{DJANGO\_SECRET\_KEY}, \texttt{ALLOWED\_HOSTS}, \texttt{CORS\_ALLOWED\_ORIGINS}, \texttt{RUN\_CODE\_ENABLED}, \texttt{GEMINI\_API\_KEY} зэрэг хувьсагчийг ашиглаж, кодоос орчны мэдээллийг салгасан.

\section{Серверийн модулиудын бүтэц}
Арын хэсгийг дараах Django модулиудаар зохион байгуулсан.
\insertimage{zurag-diagram/component-diagram.png}{Арын хэсгийн модулиудын бүрэлдэхүүний диаграм}
\begin{itemize}
\item \textbf{users}: бүртгэл, профайл, ур чадварын түвшин, эрх, нууц үгийн удирдлага
\item \textbf{courses}: суралцах замнал, курс, хичээл, таалагдсан хичээлүүд
\item \textbf{progress}: хичээл дуусгалт, оноо, ахицын хураангуй
\item \textbf{ai\_engine}: хураангуй, сорил, дадлагын даалгавар, түвшин тогтоох болон түвшин ахих логик, чатбот, код ажиллуулагч
\item \textbf{community}: пост, сэтгэгдэл, хяналт
\item \textbf{website}: нүүр хуудасны агуулга болон үнэлгээний удирдлага
\item \textbf{api}: системийн туршилтын төгсгөлийн цэг
\end{itemize}

Энэ бүтэц нь нэг модуль дотор бүх логикийг төвлөрүүлэхээс зайлсхийж, чиг үүргээр нь хуваасан тул засвар үйлчилгээ, тест, өргөтгөл хийхэд давуу талтай.

\section{Өгөгдлийн загварын хэрэгжилт}
\subsection{Хэрэглэгчийн модел}
\texttt{users.models.User} нь Django-ийн \texttt{AbstractUser}-оос удамшсан тусгай хэрэглэгчийн модел юм. Энэ моделд \texttt{role}, \texttt{display\_name}, \texttt{avatar}, \texttt{skill\_level}, \texttt{has\_placement\_test}, \texttt{total\_score}, \texttt{completed\_lessons}, \texttt{warning\_count} зэрэг талбар нэмэгдсэн. Ийм өргөтгөл нь сургалтын системийн тусгай логикийг баталгаажуулалтын системтэй шууд уялдуулах боломж олгож байна.

Display name болон avatar талбарууд нь хэрэглэгчийн интерфейсийг илүү хүн төвтэй болгоно. Ялангуяа хамтын орчны модульд хэрэглэгч хоорондын харилцаанд нэр, зураг харагдах нь оролцооны түвшинд нөлөөлдөг. \texttt{student}, \texttt{mentor}, \texttt{admin} гэсэн эрхүүд нь цаашид эрхийн илүү нарийн шатлал үүсгэх суурь болж өгнө.

\subsection{Курс ба хичээлийн модел}
\texttt{Course} модел нь \texttt{track}, \texttt{title}, \texttt{description}, \texttt{thumbnail}, \texttt{level}, \texttt{created\_by}, \texttt{created\_at} талбаруудтай. \texttt{Lesson} модел нь зөвхөн текстэн агуулгаас гадна \texttt{video\_url}, \texttt{video\_file}, \texttt{file}, \texttt{attachment}, \texttt{score}, \texttt{practice\_title}, \texttt{practice\_description}, \texttt{practice\_hint}, \texttt{practice\_expected\_output}, \texttt{practice\_bonus\_score} зэрэг талбар агуулдаг. Ийм бүтэц нь нэг хичээлийг уншлагын текст, видео тайлбар, PDF гарын авлага, хавсралт файл, дасгал, үнэлгээ гэсэн олон хэлбэрээр баяжуулах боломжтой болгож байна.

Таалагдсан хичээлийн функц нь хэрэглэгчийн сонирхлыг илэрхийлэх анхны зан төлөвийн дохио болдог. Цаашид зөвлөмжийн логикийг өргөтгөх үед таалагдсан хичээлүүд, дахин үзсэн байдал, дуусгалтын хурд зэрэг мэдээлэл ашиглах боломжтой.

\subsection{Ахицын модел}
\texttt{LessonProgress} нь хэрэглэгч ба хичээл хоорондын холбоос бөгөөд \texttt{is\_completed}, \texttt{score}, \texttt{completed\_at} мэдээллийг хадгална. \texttt{unique\_together} нөхцөлөөр нэг хэрэглэгч нэг хичээл дээр нэг progress мөртэй байхаар хамгаалсан. Энэхүү модел нь ахицын тооцооллын үндсэн эх сурвалж юм.

Суралцагч нэг хичээлийг үзсэн, дуусгасан, өндөр оноотой дуусгасан гэх мэт олон төлөвтэй байж болно. Энэхүү төсөлд \texttt{LessonProgress} нь эдгээрээс хамгийн чухал гурван мэдээллийг хадгалдаг. Мөн \texttt{complete\_lesson} төгсгөлийн цэгт score утгыг \texttt{clamp\_score} функцээр дээд, доод хязгаарт барьдаг тул урд талын хэсгээс буруу утга ирсэн ч логик бус өгөгдөл хадгалагдах эрсдэл буурна.

\subsection{Community моделийн бүтэц}
\texttt{CommunityPost} болон \texttt{CommunityComment} моделууд нь \texttt{author}, \texttt{content}, \texttt{moderation\_status}, \texttt{moderation\_reason}, \texttt{created\_at} мэдээллийг хадгална. Пост нь хичээлтэй сонголтот холбоотой тул хичээл бүр дээр хэлэлцүүлгийн орчин үүсгэх боломж бүрдсэн. \texttt{warning\_count} талбар нь хяналтын шийдвэрийг хэрэглэгчийн урт хугацааны зан төлөвтэй холбох эхний алхам болж байна.

\section{Нэвтрэлт ба эрхийн удирдлага}
Баталгаажуулалтын хэсэгт JWT ашигласан. Хэрэглэгч нэвтрэх үед access token үүсч, урд талын хэсэг энэ token-ийг хадгалан дараагийн API хүсэлтүүдэд ашиглана. DRF-ийн \texttt{IsAuthenticated} эрхийг курс, хичээл, ахиц, хамтын орчин, хиймэл оюуны хураангуй зэрэг хамгаалагдсан төгсгөлийн цэгүүдэд өргөн ашигласан.

Эрхийн удирдлагын хувьд \texttt{IsStaffMentorOrReadOnly} гэсэн тусгай эрх ашигласан. Энэ нь зөвхөн унших үйлдлийг энгийн нэвтэрсэн хэрэглэгчдэд нээлттэй байлгаж, үүсгэх болон шинэчлэх төрлийн үйлдлийг admin, staff, mentor хэрэглэгчдээр хязгаарлахад зориулагдсан.

\section{Урд талын архитектур}
Урд талын чиглүүлэлт нь \texttt{BrowserRouter}, \texttt{Routes}, \texttt{ProtectedRoute}, \texttt{DashboardLayout}, \texttt{AuthLayout} зэрэг бүрэлдэхүүнээр зохион байгуулагдсан. Үндсэн хуудаснууд дараах байдалтай байна.
\begin{itemize}
\item Нүүр танилцуулгын хуудас
\item Нэвтрэх, бүртгүүлэх хуудас
\item Хяналтын самбарын нүүр
\item Курсийн жагсаалт, курсийн дэлгэрэнгүй
\item Хичээлийн дэлгэрэнгүй
\item Замналууд, замналын курсүүд
\item Миний суралцах хэсэг, ахицын хураангуй
\item Түвшин тогтоох шалгалт, түвшин ахих шалгалт
\item Хамтын орчин
\item Тоглоомын хэсэг, тоглоомын дэлгэрэнгүй
\item Хувийн мэдээлэл
\item Таалагдсан хичээлүүд
\end{itemize}

\texttt{ProtectedRoute} нь нэвтрээгүй хэрэглэгчийг хамгаалагдсан чиглэлүүдээс сэргийлнэ. \texttt{DashboardLayout} нь хажуугийн цэс, дээд мөр, үндсэн контентын хүрээг нэг мөр хэв загвартай болгодог. \texttt{PageFade} болон \texttt{framer-motion} ашигласнаар хуудас хоорондын шилжилт илүү зөөлөн болсон.

\section{Сургалтын үндсэн хуудсуудын хэрэгжилт}
\subsection{Нүүр танилцуулгын хуудас ба хяналтын самбар}
Нүүр танилцуулгын хуудас нь системийн зорилго, онцлог, хэрэглэгчийн сэтгэгдэл, бүртгүүлэх уриалгыг харуулдаг. Энэхүү хуудасны агуулгыг код дотор шууд бичээгүй, харин \texttt{LandingPageContent} болон \texttt{LandingPageReview} моделуудаар өгөгдлийн сангаас удирддаг. Хяналтын самбарын нүүр нь хэрэглэгч нэвтэрсний дараах үндсэн навигацийн төв цэг бөгөөд түвшинд тохирсон агуулга, онцлох курс, товчилсон санал болголтыг харуулахад зориулагдсан.

\subsection{Курс, замнал ба таалагдсан хичээлүүд}
Курсийн хуудас нь агуулгыг түвшин, замналын дагуу хайх, шүүх орчин болдог. Курсийн дэлгэрэнгүй нь тухайн курсэд харьяалагдах хичээлүүдийг дарааллаар харуулж, суралцах дараагийн алхмыг ойлгомжтой болгоно. Замналын хэсэг нь ижил чиглэлийн курсүүдийг бүлэглэж, хэрэглэгчийг навигацид төөрөхөөс сэргийлнэ. Таалагдсан хичээлүүд нь хэрэглэгчийн сонирхсон хичээлүүдийг дахин хурдан нээхэд зориулагдсан.

\subsection{Хичээлийн дэлгэрэнгүй хуудас}
Хичээлийн дэлгэрэнгүй хуудас нь системийн хэрэглээний гол цөм юм. Энэ хуудсанд дараах боломжууд нэг дор төвлөрдөг.
\begin{itemize}
\item Хичээлийн текстэн агуулга үзэх
\item Видео файл эсвэл URL үзэх
\item Хиймэл оюуны хураангуй гаргах
\item Хиймэл оюуны сорил ажиллуулах
\item Дадлагын даалгавар унших, хариулт бичих
\item Веб урьдчилан харах эсвэл Python код ажиллуулах
\item Хамтын орчны хэлэлцүүлэг үзэх, бичих
\item Хичээлд таалагдсан тэмдэглэгээ хийх
\item Хичээл дууссаныг баталгаажуулах
\end{itemize}

Энэхүү зохион байгуулалт нь хэрэглэгчийг олон тусдаа хуудас дамжуулахгүйгээр суралцах үйл явцыг тасралтгүй барих давуу талтай.

\subsection{Түвшин тогтоох шалгалт ба түвшин ахих шалгалт}
Түвшин тогтоох шалгалт нь хэрэглэгчийн анхны түвшнийг тогтоох шийдвэрийн цэг юм. Түвшин ахих шалгалт нь одоогийн түвшинг давсан эсэхийг шалгаж, дараагийн шат руу ахиулахад ашиглагдана. Эдгээр нь ердийн сорилын хуудас биш, харин системийн сургалтын урсгалыг өөрчилдөг механизм учраас бүтцийн хувьд тусгай байр суурь эзэлдэг.

\subsection{Ахицын хураангуй, хамтын орчин ба тоглоомын хэсэг}
Ахицын хураангуй нь нийт оноо, дуусгасан хичээл, долоо хоногийн идэвх, курс тус бүрийн гүйцэтгэлийг нэгтгэн харуулдаг шинжилгээний хуудас юм. Хамтын орчин нь хэрэглэгч хоорондын харилцааг, тоглоомын хэсэг нь идэвхжүүлэлтийг дэмждэг. Иймд систем нь зөвхөн агуулга үзэх статик платформ бус, харилцаа болон оролцоотой сургалтын орчин болж өргөжсөн.

\section{Шаталсан сургалтын логикийн хэрэгжилт}
\subsection{Түвшин тогтоох шалгалт}
Түвшин тогтоох шалгалт нь хэрэглэгчийн эхний түвшнийг тогтооно. \texttt{ai\_engine} дотор \texttt{collect\_level\_source\_text("beginner")} функцээр beginner түвшний курс болон хичээлүүдийн текстийг цуглуулж, үүн дээр үндэслэн сорил үүсгэх бүтэц бэлтгэсэн. Хэрэглэгчийн хариултыг илгээсний дараа зөв хариултын хувь тооцоологдон, \texttt{save\_level} төгсгөлийн цэгээр \texttt{skill\_level} болон \texttt{has\_placement\_test} талбар шинэчлэгдэнэ.

\subsection{Түвшин ахих шалгалт}
Түвшин ахих шалгалт нь хэрэглэгч дараагийн шатанд шилжихэд хэрэглэгдэнэ. Одоогийн логик нь хэрэглэгчийн одоогийн түвшинд хамаарах эх текстийг авч, түүнээс шат ахих шалгалтын асуултууд үүсгэдэг. Энэ нь тогтмол асуултын сантай байхгүйгээр системийг уян хатан байлгадаг.

\subsection{Санал болгох логик}
\texttt{recommended\_lessons} төгсгөлийн цэг нь хэрэглэгчийн дуусгасан хичээлүүдийг хасаж, одоогийн түвшинд тохирох хичээлүүдээс эхний тавыг санал болгодог. Энэ нь дүрэмд суурилсан боловч эхний хувилбарын хувьд тайлбарлахад хялбар, хэрэгжүүлэхэд ойлгомжтой шийдэл болсон.

\section{Хиймэл оюуны хөдөлгүүрийн хэрэгжилт}
\insertimage{zurag/diagram_ai.png}{Хиймэл оюуны хөдөлгүүрийн модулийн ерөнхий урсгал}

\subsection{Промпт боловсруулах зарчим}
Хиймэл оюуны endpoint бүрт prompt-ийг тодорхой дүрэмтэйгээр өгсөн. Жишээлбэл quiz generation-д JSON-оос өөр текст бичихгүй байх, дөрвөн сонголттой байх, \texttt{answer\_index} 0-оос эхлэх, тайлбар заавал өгөх зэрэг хязгаарлалт тавьсан. Энэ нь LLM-ийн хариуг backend дээр parse хийх боломжтой байлгах гол нөхцөл юм.

\subsection{Хураангуй үүсгэх логик}
Хичээлийн хураангуйн төгсгөлийн цэг нь эхлээд хичээлийн \texttt{content} талбарыг авна. Хэрэв агуулга хоосон бөгөөд зөвхөн PDF файл байвал текст ялгалт хийгдээгүй тухай мэдэгдэл буцаана. Харин агуулга байгаа тохиолдолд \texttt{llm\_summarize} функцийг дуудаж Gemini API-аар хураангуй үүсгэнэ.
\insertimage{zurag-diagram/ai-summary.png}{Хичээлийн хураангуй үүсгэх хиймэл оюуны урсгалын дарааллын диаграм}

Хэрэв API түлхүүр байхгүй эсвэл хиймэл оюуны хүсэлт алдаа өгвөл \texttt{simple\_summarize} нөөц логик ажиллана. Энэ нөөц логик нь текстийг өгүүлбэрээр таслан эхний хэсгүүдийг авдаг энгийн арга боловч системийн тасралтгүй ажиллагааг хангах чухал үүрэгтэй.

\subsection{Сорил үүсгэх логик}
\texttt{llm\_quiz} болон \texttt{build\_quiz\_from\_source\_text} функцууд нь хиймэл оюунаар multiple-choice асуулт үүсгэдэг. Хариу JSON формат алдах магадлалтай тул backend parse шатанд \texttt{try/except} ашиглан бүтэц алдагдсан асуултуудыг шүүж, зөв форматтай хэсгийг л frontend-д дамжуулдаг.

\subsection{Дадлагын даалгавар үүсгэх логик}
Дадлагын даалгавар байхгүй хичээлд \texttt{ensure\_practice\_task\_for\_lesson} функц ажиллаж, title, description, hint, expected output бүхий даалгавар үүсгэн хадгална. Ингэснээр дараагийн удаа ижил хичээл нээгдэхэд хиймэл оюуныг дахин заавал дуудах шаардлагагүй болно. Хэрэв хиймэл оюун ашиглах боломжгүй бол анхны жижиг даалгаврын бүтэц буцаадаг.

\subsection{Код ажиллуулах хэсэг ба sandbox}
Код ажиллуулагчийн хэсэгт \texttt{apply\_runner\_limits} функцээр CPU хугацаа, санах ой, файлын гаралтын хэмжээнд хязгаар тогтоосон. Энэ нь суралцагчийн кодыг ажиллуулах үед системийн нөөцийг хамгаалах, хяналтгүй урт процесс эсвэл хэт их гаралт үүсэх эрсдэлийг бууруулах зорилготой.

\section{Ахицын тооцоолол ба хэрэглэгчийн төлөв шинэчлэлт}
\subsection{Хичээл дуусгалт}
\texttt{complete\_lesson} endpoint нь хичээлийг дууссанд тооцох бөгөөд score параметр ирсэн эсэхээс хамааран quiz эсвэл default lesson score ашиглана. Тухайн progress мөрийг хадгалсны дараа хэрэглэгчийн дуусгасан хичээлийн тоо болон нийт оноог нэгтгэн тооцоолж шинэчилдэг.

\subsection{Хэрэглэгчийн түвшин шинэчлэлт}
\texttt{User.update\_level\_and\_role()} функц нь нийт оноонд тулгуурлан хэрэглэгчийн түвшинг шинэчилнэ. Хэрэв хэрэглэгч advanced түвшинд хүрч, дор хаяж таван хичээл дуусгасан бол mentor role руу шилжүүлдэг.

\subsection{Долоо хоногийн идэвх ба ахицын хураангуй}
\texttt{progress\_summary} endpoint нь хэрэглэгчийн \texttt{username}, \texttt{role}, \texttt{skill\_level}, \texttt{total\_score}, \texttt{completed\_lessons}, долоо хоногийн идэвх, курс бүрийн ахицын хувь, курсын оноо зэрэг мэдээллийг нэгтгэн буцаана. Долоо хоногийн идэвхийг \texttt{TruncDate} ашиглан өдрөөр бүлэглэж гаргасан нь шинжилгээний эхний хувилбар болж байна.

\section{Хамтын орчны хяналтын хэрэгжилт}
Хамтын орчны хэсэг нь зөвхөн нийтлэл бичих энгийн форум биш, хичээлтэй уялдсан хэлэлцүүлгийн орчин байдлаар төлөвлөгдсөн. Пост болон сэтгэгдэл үүсгэх үед \texttt{moderate\_text} функц ажиллаж эхлээд түлхүүр үгийн хяналт, дараа нь хиймэл оюуны хяналт хийнэ.

Түлхүүр үгийн хяналт нь хямд, хурдан, тогтвортой механизм юм. Харин хиймэл оюуны хяналт нь илүү уян хатан боловч алдаа гарах магадлалтай. Иймээс \texttt{allow}, \texttt{review}, \texttt{reject} гэсэн гурван төлөвийг хослуулсан шийдэл ашигласан. Хэрэв хориглосон үг илэрвэл \texttt{reject} шийдвэр гарч, хэрэглэгчийн \texttt{warning\_count} нэмэгдэнэ.

\section{API урсгалын жишээ}
\subsection{Хичээл дуусгалтын API}
\insertimage{zurag-diagram/lesson-comp.png}{Хичээл дуусгалтын API-ийн дарааллын диаграм}
\begin{enumerate}
\item Урд талын хэсэг хэрэглэгчийн авсан score эсвэл анхны score-г арын хэсэг рүү илгээнэ.
\item Арын хэсэг хичээл байгаа эсэхийг шалгана.
\item User болон lesson хослол дээр progress бичлэг үүсгэх эсвэл сэргээн авна.
\item Progress дээр completion төлөв, score, timestamp хадгална.
\item Хэрэглэгчийн дуусгасан хичээлийн тоо, нийт оноог дахин тооцоолно.
\item Ур чадварын түвшин болон эрх шинэчилнэ.
\item JSON хариу буцаана.
\end{enumerate}

\subsection{Хичээлийн хураангуйн API}
\begin{enumerate}
\item Урд талын хэсэг хураангуйн хүсэлт илгээнэ.
\item Арын хэсэг хичээлийн агуулгыг авна.
\item Хэрэв агуулга хоосон бөгөөд зөвхөн PDF байвал мэдээллийн мэдэгдэл буцаана.
\item Эс бөгөөс хиймэл оюуны хураангуй үүсгэхийг оролдоно.
\item Хиймэл оюун амжилтгүй бол нөөц хураангуй ажиллана.
\item Хураангуй JSON хариу хэлбэрээр урд талын хэсэгт хүрнэ.
\end{enumerate}

\section{Docker орчин ба нэвтрүүлэлтийн бэлтгэл}
\texttt{docker-compose.yml} файлд \texttt{db} болон \texttt{backend} гэсэн хоёр үндсэн үйлчилгээ тодорхойлсон. PostgreSQL үйлчилгээ нь бэлэн байдлын шалгалттай бөгөөд \texttt{backend} нь өгөгдлийн сан бэлэн болмогц асдаг. \texttt{wait-for-db.sh} скрипт ашигласнаар migration эсвэл runserver эхлэх үед өгөгдлийн сангийн бэлэн байдлын асуудлыг бууруулсан.

Орчны хувьсагч ашигласнаар локал болон үйлдвэрлэлийн орчны ялгааг уян хатан удирдах боломжтой. Үйлдвэрлэлийн шалгах жагсаалт дотор \texttt{DJANGO\_DEBUG=0}, secure cookie, SSL redirect, HSTS, бодит домэйн, \texttt{VITE\_API\_BASE\_URL} зэрэг шаардлагуудыг тусгасан.

\section{Аюулгүй байдал ба өргөтгөх боломж}
Системд JWT баталгаажуулалт, CORS allowed origins-ийн нарийн тохиргоо, CSRF trusted origins, secure cookie, HTTPS redirect, HSTS, код ажиллуулагчийн CPU ба санах ойн хязгаар, хамтын орчны хяналт зэрэг арга хэмжээг тусгасан. Эдгээр нь бүрэн хэмжээний үйлдвэрлэлийн аюулгүй байдлын аудитийг орлохгүй ч төслийн түвшинд бодит эрсдэлийг инженерчлэлийн аргаар тооцсон шийдэл болсон.

Одоогийн зохиомж нь admin болон mentor-д зориулсан курсийн агуулга удирдах самбар, зөвлөмжийн хөдөлгүүрийг машин сургалт эсвэл хосолмол алгоритмаар сайжруулах, PDF текст ялгалт, шинжилгээний самбар, гар утсанд чиглэсэн хувилбар, хамтын орчны нэр хүндийн систем зэрэг өргөтгөлүүдэд бэлэн байна.

\section{Бүлгийн дүгнэлт}
Энэ бүлэгт системийн архитектур, технологийн стек, арын хэсгийн модуль, өгөгдлийн загвар, баталгаажуулалт, урд талын бүтэц, хичээлтэй холбоотой төв логик, хиймэл оюуны хураангуй, сорил, дадлагын даалгавар үүсгэх логик, ахиц хянах, хяналт, Docker орчин, аюулгүй байдлын шийдлийг нэгтгэн тайлбарлав. Ийнхүү өмнө нь тусдаа авч үзэж байсан модуль тус бүрийн нарийвчилсан тайлбаруудыг системийн зохиомжийн хүрээнд шингээснээр энэхүү бүлэг хөгжүүлсэн платформын техникийн хэрэгжилтийг бүрэн баримтжуулсан хэсэг боллоо.
